\documentclass[12pt]{article}
\usepackage{amsmath, amssymb}

\title{Lecture 21--22: Inequalities and Tail Bounds}
\author{CSE400 -- Fundamentals of Probability in Computing}
\date{}

\begin{document}

\maketitle

\section*{1. Introduction and Motivation}

\subsection*{1.1 What is a Tail?}

\textbf{Definition:} \\
The \textit{tail} of a random variable $(X)$ is:
\[
P(X > x)
\]

This represents the probability that the random variable takes \textbf{large values beyond a threshold $(x)$}. 

\subsection*{1.2 Why Do We Care About Tails?}

From the slides (page 5):

\begin{itemize}
\item Jobs delayed more than 24 hours
\item Packets dropped due to full buffer
\end{itemize}

These represent \textbf{rare but critical events}.

\subsection*{1.3 Key Idea}

\begin{itemize}
\item \textbf{Mean $\rightarrow$ average behavior (easy to compute)}
\item \textbf{Tail $\rightarrow$ rare extreme events (hard to compute)} 
\end{itemize}

Thus, while averages describe normal behavior, \textbf{tails capture worst-case scenarios}, which are crucial in systems.

\subsection*{1.4 Why Are Tails Hard to Compute?}

\begin{itemize}
\item Requires \textbf{summing many small probabilities}
\item Leads to \textbf{long and complex expressions} 
\end{itemize}

\section*{2. Examples of Tail Events}

\subsection*{2.1 Example 1: Load Balancing (Binomial Setting)}

\textbf{Problem:} \\
Distribute $(n)$ jobs among $(n)$ servers randomly. \\
Find:
\[
P(\text{a server gets } \ge k \text{ jobs})
\]

\subsubsection*{Intuition (from slides):}

\begin{itemize}
\item Jobs are assigned randomly
\item Most servers get few jobs
\item Occasionally, one server gets many jobs
\end{itemize}

This is a \textbf{tail event (rare overload)} 

\subsubsection*{Mathematical Model:}

\[
X \sim \text{Binomial}(n, p)
\]

\[
P(X \ge k) = \sum_{i=k}^{n} \binom{n}{i} p^i (1-p)^{n-i}
\]

\subsection*{2.2 Example 2: Poisson Arrivals}

\textbf{Problem:} \\
Jobs arrive at rate $(\lambda)$ (Poisson process). \\
Find:
\[
P(\text{at least } k \text{ arrivals in 1 hour})
\]

\subsubsection*{Intuition:}

\begin{itemize}
\item Jobs arrive randomly over time
\item Most hours $\rightarrow$ moderate arrivals
\item Sometimes $\rightarrow$ burst of arrivals
\end{itemize}

This is a \textbf{tail event (sudden spike)} 

\subsubsection*{Mathematical Model:}

\[
X \sim \text{Poisson}(\lambda)
\]

\[
P(X \ge k) = \sum_{i=k}^{\infty} e^{-\lambda} \frac{\lambda^i}{i!}
\]

\section*{3. Why Tail Bounds Are Needed}

\subsection*{3.1 Exact Computation is Difficult}

From slides (page 14):

\begin{itemize}
\item Binomial and Poisson tails require \textbf{long summations}
\item These are:
\begin{itemize}
\item Computationally expensive
\item Hard to simplify
\end{itemize}
\end{itemize}

\subsection*{3.2 Key Idea: Approximate Instead of Exact}

Instead of computing:
\[
P(X \ge k)
\]

We compute an \textbf{upper bound}:
\[
P(X \ge k) \le \text{simple expression}
\]

\subsection*{3.3 Advantages}

\begin{itemize}
\item Easier to compute
\item Still provides a \textbf{guarantee}
\item Useful in system design
\end{itemize}

\subsection*{3.4 Core Philosophy}

\begin{quote}
\textbf{``Approximate safely instead of computing exactly.''}
\end{quote}

\section*{4. Running Example (Used Across Inequalities)}

We analyze all bounds using the same setup:

\subsection*{Experiment:}

\begin{itemize}
\item Flip a fair coin $(n)$ times
\end{itemize}

\subsection*{Random Variable:}

\[
X = \text{number of heads}
\]

\[
X \sim \text{Binomial}(n, \tfrac{1}{2})
\]

\subsection*{Mean:}

\[
E[X] = \frac{n}{2}
\]

\subsection*{Tail Question:}

\[
P\left(X \ge \frac{3n}{4}\right)
\]

This represents \textbf{getting unusually many heads}, i.e., a \textbf{tail region}. 

\section*{5. Markov's Inequality}

\subsection*{5.1 Key Idea}

\begin{itemize}
\item Uses \textbf{only the mean}
\item Large values are \textbf{rare}
\end{itemize}

Example intuition:

\begin{itemize}
\item If average = 10 $\rightarrow$ very few values can be $\ge 100$ 
\end{itemize}

\subsection*{5.2 Statement of Markov's Inequality}

For any non-negative random variable $(X)$ and $(a > 0)$:

\[
P(X \ge a) \le \frac{E[X]}{a}
\]

\subsection*{5.3 Interpretation}

\begin{itemize}
\item Probability of large deviation is bounded using \textbf{mean only}
\item No need for full distribution
\end{itemize}

\subsection*{5.4 Logical Flow}

\begin{enumerate}
\item We want $(P(X \ge a))$
\item Exact computation is hard
\item Use only $(E[X])$
\item Obtain a simple upper bound
\end{enumerate}

\section*{6. Overall Conceptual Flow (Important for Exams)}

\begin{enumerate}
\item \textbf{Tail definition:} \\
$(P(X > x))$

\item \textbf{Motivation:} \\
Rare events (failures, overloads)

\item \textbf{Problem:} \\
Exact computation requires large sums

\item \textbf{Solution:} \\
Use \textbf{tail bounds}

\item \textbf{First tool:} \\
Markov's Inequality

\item \textbf{Key principle:} \\
Use simple statistics (like mean) to bound complex probabilities
\end{enumerate}

\section*{7. Key Takeaways}

\begin{itemize}
\item Tail probabilities measure \textbf{extreme events}
\item Exact computation is \textbf{difficult due to summations}
\item Tail bounds provide:
\begin{itemize}
\item Simplicity
\item Guarantees
\end{itemize}
\item Markov inequality is the \textbf{simplest bound}
\item More advanced bounds (Chebyshev, Chernoff) improve tightness (as per outline)
\end{itemize}

\end{document}